\subsection{Metavariables and higher-order pattern unification}\label{sec:miller}

Tactics and the elaborator work with partially built terms.
A goal is a metavariable $\meta m$ with a type and the context in which it was created; \code{apply} introduces a metavariable for each argument of a lemma that the goal does not determine; \code{exact} and \code{rewrite} leave constraints for terms that are not yet known.
Resolving such constraints needs higher-order unification, which is undecidable in general \citep{huet75}.
The kernel restricts it, as every production elaborator does, to the fragment of Miller's \emph{patterns} \citep{miller91,nipkow93}, which has unique most general solutions, and defers everything else.

\paragraph{State.}
Metavariables live in a metavariable context: each has a type and, once assigned, a solution, a closed term.
Assignment is recorded on a trail, so a context can be checkpointed and rewound in time proportional to what was assigned since.
The elaboration state also holds a list of \emph{pending} constraints, each a pair of terms with the context in which it was deferred.
A term in which a solved metavariable occurs is read through the solution: a solved metavariable at the head of a neutral term is replaced by its solution before the term is compared.

\begin{definition}[Pattern]\label{def:pattern}
A neutral term is a \emph{Miller pattern} if its head is a metavariable $\meta m$ and its spine consists of applications to \emph{distinct bound variables of the context}, each a bare variable and not itself applied.
In particular a pattern has no case-analysis frame in its spine.
\end{definition}

For a pattern $\meta m\,x_1\cdots x_k$ and a term $t$ with $\meta m\notin t$ and $\FV(t)\subseteq\{x_1,\dots,x_k\}$, the term $\lambda x_1\cdots x_k.\,t$ is the most general solution of $\meta m\,x_1\cdots x_k \conv t$ \citep{miller91}.
Because $\meta m$ is replaced by a closed abstraction, the solution is well scoped wherever $\meta m$ occurs.
In a dependent calculus the telescope of binders needs types, and \textsc{Solve} (\cref{alg:solve}) takes them from the types of the spine variables, with two additional checks for the dependency that the simply typed setting does not have.

\begin{algorithm}[t]
\caption{\Solve$(m,\vs x,t)$: assign the pattern $\meta m\,\vs x \conv t$ in a context $\Gamma$ of size $n$}
\label{alg:solve}
\begin{algorithmic}[1]
\Require a metavariable $m$ unassigned; distinct context variables $\vs x = x_{l_1},\dots,x_{l_k}$ with types $D_1,\dots,D_k$ taken from $\Gamma$; a term $t$ (read back from a value)
\Ensure the assignment of $m$, or a failure
\If{$m$ occurs in $t$} \Return failure \Comment{occurs check}
\EndIf
\State $S \gets \FV(t)\ \cup\ \bigcup_{i}\FV(D_i)$
\If{$S\not\subseteq\{x_{l_1},\dots,x_{l_k}\}$} \Return failure \Comment{scope: the solution may depend only on the spine, and so may the types of the spine}
\EndIf
\State let $y_1,\dots,y_k$ be fresh context variables beyond $\Gamma$, $\varrho = [x_{l_i}\mapsto y_i]$
\For{$i = 1,\dots,k$}
   \State $D_i' \gets \varrho(D_i)$
   \If{$D_i'$ mentions $y_i,\dots,y_k$} \Return failure \Comment{the spine must be in dependency order}
   \EndIf
\EndFor
\State assign $m \gets \lambda(y_1{:}D_1')\cdots(y_k{:}D_k').\ \varrho(t)$ \Comment{closed over the fresh telescope}
\State \Return success
\end{algorithmic}
\end{algorithm}

\begin{algorithm}[t]
\caption{\Unify$(a,b)$ in a context $\Gamma$ of size $n$, on values; the pending list and the metavariable context are state}
\label{alg:unify}
\begin{algorithmic}[1]
\State $a \gets \sy{resolve}(a)$, $b \gets \sy{resolve}(b)$ \Comment{replace a solved metavariable at the head by its solution}
\If{$a$ and $b$ are meta-headed and $\qu_n(a) = \qu_n(b)$} \Return success \Comment{$\meta m\,\vs x \conv \meta m\,\vs x$}
\EndIf
\If{$a$ is a pattern $\meta m\,\vs x$} \Return \Solve$(m,\vs x,\qu_n(b))$ \Comment{the left side is tried first}
\EndIf
\If{$b$ is a pattern $\meta m\,\vs x$} \Return \Solve$(m,\vs x,\qu_n(a))$
\EndIf
\If{$a$ or $b$ is meta-headed} \State push $(\qu_n(a),\qu_n(b),\Gamma)$ onto \textit{pending}; \Return success \Comment{deferred, not decided}
\EndIf
\State compare $a$ and $b$ by shape:
\Statex \hspace{1.5em}\textbullet\ two sorts: equal levels;
\Statex \hspace{1.5em}\textbullet\ two closures, or two products: \Unify\ the domains, then the bodies applied to a fresh neutral term, in $\Gamma$ extended by the domain;
\Statex \hspace{1.5em}\textbullet\ a closure against a neutral term or a fixpoint: apply the other side to a fresh neutral term ($\eta$) and \Unify\ the bodies;
\Statex \hspace{1.5em}\textbullet\ two constructor values: same constructor, \Unify\ the fields pairwise;
\Statex \hspace{1.5em}\textbullet\ two neutral terms: same head and spine length, \Unify\ application frames pairwise; stuck case frames must have the same motive and branches convertible in their environments; on an uncertain failure compare read-backs;
\Statex \hspace{1.5em}\textbullet\ anything else: failure.
\If{the comparison failed} \State roll back the metavariables and pending constraints it touched; try \textsc{PI} (\cref{sec:prop}) \EndIf
\end{algorithmic}
\end{algorithm}

\paragraph{Unification.}
\Cref{alg:unify} is the unifier.
It extends the conversion of \cref{sec:conv} with the two pattern cases.
Two solved-or-unsolved applications of the same metavariable to the same arguments are equal (the reference kernel of \cref{sec:testing} failed on $\meta m\,\vs x\conv\meta m\,\vs x$ because its occurs check fired on the problem itself).
When both sides are patterns the left one is solved.
A meta-headed problem that is not a pattern ($\meta m\,(g\,x)\conv c$, or $\meta m\,x\,x\conv c$ with a repeated variable) is neither solved nor failed: it is put on the pending list.

\paragraph{Draining.}
After a unification that may have solved a metavariable, the elaborator re-attempts every pending constraint in the context in which it was deferred, and repeats until a pass makes no progress (the number of pending constraints does not decrease) or the list is empty.
A constraint that has become a pattern, or whose heads have been solved, is then decided.
\emph{Acceptance} of an elaboration requires that the pending list is empty: a tactic whose proof state is not fully resolved cannot be accepted, and a term whose constraints are not all decided cannot be admitted.
This gate is not negotiable, because a deferred constraint is a promise that something will be compared, and a state that holds a promise it never keeps looks as if it had been checked.

\paragraph{What the algorithm is, and is not.}
\Cref{alg:unify} is the pattern fragment of the textbook algorithm with the following departures, which make it weaker than a full implementation; none of them can cause an ill-typed term to be admitted, since the final term is checked in full.
\begin{itemize}
\item \emph{No pruning.} A metavariable $\meta n$ inside $t$ that is applied to variables outside the spine of $\meta m$ is not restricted to the variables of the spine: the scope check fails instead.
  Systems with pruning \citep{abel11,ziliani15} solve more problems.
\item \emph{No type check of the solution.} \Solve\ checks scope and dependency order and does not check that the solution has the type of $\meta m$.
  The solution is type-correct whenever the problem was, because metavariables are created with the type of their context and every unification is of terms of one type; but the unifier does not prove this, and the elaboration of \code{apply} re-checks each solved argument metavariable on a scratch copy (finding F25 in \cref{sec:testing}).
  What is admitted is the \emph{final} term, which is checked in full.
\item \emph{Patterns by construction.} The metavariable of a goal is applied to every variable of its context, hence is a pattern; the restriction to patterns bites only on the metavariables of implicit arguments.
\item \emph{Deferral instead of search.} There is no backtracking over alternative solutions (no search for imitation and projection instances \citep{huet75}); a problem that is not a pattern waits for other assignments or remains pending and refuses acceptance.
\end{itemize}

\begin{example}
In a context $x{:}\Nat,\,y{:}\Nat$ the problem $\meta m\,x\,y \conv \Fam{add}\,y\,x$ is a pattern problem: $S=\{x,y\}$, the dependency order is trivial, and the assignment is $\meta m := \lam{x}{\Nat}{\lam{y}{\Nat}{\Fam{add}\,y\,x}}$.
The problem $\meta m\,(\ctor{Succ}\,x)\conv c$ is deferred.
The problem $\meta m\,x\,x \conv c$, with a repeated variable, is deferred as well, and is decided only if $\meta m$ is assigned by another constraint.
\end{example}

\paragraph{Termination and backtracking.}
We do not argue termination of unification from a measure; it is enforced by the work budget of \cref{sec:impl}, which makes every call end in success, in failure, or in a budget error that fails closed.
A metavariable is assigned at most once between a checkpoint and the rewind that follows it, and every tactic runs in a transaction that rewinds the metavariable context and the pending list on any exception, so that a failed attempt leaves no assignment behind.

\subsection{The trust boundary}\label{sec:trust}

Admission to the global environment is the only operation that changes what is true.
\begin{definition}[Checked entry points]\label{def:entry}
The global environment is extended by $\Sg \to \Sg,\,c : A \mathrel{:=} v$ only when
\[
\inferrule*[right=\pf{Declare}]{c\notin\dom(\Sg) \\ \Sg;\cdot\vdash A : \Type_i \\ \Sg;\cdot\vdash v : A \\ A, v \text{ contain no metavariable and no unfinished global}}
{\Sg \;\longrightarrow\; \Sg,\ c : A \mathrel{:=} v}
\]
and by the three other checked operations: \texttt{finish\_inductive} (admission conditions \textbf{A1}--\textbf{A6}), \texttt{finish\_recursive\_function\_checked} (the body checked against its type under the guard), and \texttt{declare\_quot} (the preconditions of \cref{sec:quot}).
\end{definition}
The check of \textsc{Declare} is made in the \emph{empty} context against a closed statement.
Everything on the other side of this door (the tactic engine, the decision procedures, the elaborator and its unifier, the environment, the search procedures and the learner) is untrusted: it constructs candidate terms, and its only power over the environment is to submit them.
A bug in any of it can fail to find a proof, or make a candidate that the kernel then refuses; it cannot extend the environment with an unchecked declaration.

The unchecked variants (\texttt{declare\_unchecked}, \texttt{finish\_inductive\_unchecked}) exist, for axioms and so that the tests can build ill-formed environments on purpose.
In the code we ship, the only call to one of them is the declaration of an axiom, and the libraries declare exactly one axiom, excluded middle, in the number-theory library.
The function \texttt{audit} reports, for a theorem, the axioms and quotient constants that it reaches through the types and values of the definitions it uses, in the manner of Lean's \texttt{\#print axioms}; a theorem that reaches none depends only on declarations that entered by a checked entry point, provided the discipline holds.
The module that exposes the kernel to Python does expose the unchecked variants; the boundary is therefore a discipline of the layers above the kernel, backed by the second checker (\cref{sec:indep}), which re-checks the value of every definition in the exported environment, whichever path installed it.
It does not re-check the admission of inductive families.

\subsection{Certificate calculi}\label{sec:certificates}

The decision procedures are \emph{certifying}: each produces a proof term, built from lemmas of the libraries, that the kernel checks like any other.
They are untrusted in the sense of \cref{sec:trust}: a wrong search produces a failure or a refused term, never a theorem.
Each is described by the certificate that it searches for and the lemma chain that turns it into a term.

\paragraph{Ring normalisation.}
The tactic \code{ring} decides equality in the commutative semirings of the unary naturals $\Fam{Nat}$ and the binary naturals $\Fam{N}$ and in the commutative ring of the integers $\Fam{Z}$ of the libraries, by normalising both sides to a canonical polynomial and building, at each step, the application of a ring lemma that justifies it \citep{gregoire05,mathlib20}.
The normal form follows Mathlib's:
\[
\text{poly} \;::=\; m_1 + (m_2 + (\cdots + (m_k + 0)))
\qquad
\text{mono} \;::=\; x_1^{e_1}\cdot(x_2^{e_2}\cdot(\cdots\cdot(x_j^{e_j}\cdot c)))
\]
with monomials in strictly decreasing order, atoms in increasing index order, exponents $e\ge1$ numerals and the coefficient $c$ a numeral last.
The order on monomials is the lexicographic order on exponent vectors, with the first atom most significant; it is total and compatible with multiplication, which is what lets a product of a monomial and a polynomial stay sorted without merging.
Numerals are the canonical numerals of the number system, and arithmetic on them is checked by the kernel's evaluation (in time proportional to the number of bits for binary numerals).
A power by a numeral uses repeated squaring, so $(a+b)^{20}$ takes five multiplications.
In a ring, subtraction is multiplication by the numeral $-1$, and monomials whose coefficients cancel leave the sum.
If the normal forms differ, the tactic fails and reports both.

\paragraph{Linear arithmetic by Farkas certificates.}
The tactic \code{linarith} reads each hypothesis as a fact $0\le p_i$ with $p_i$ a polynomial over $\Z$ ($x\le y$ becomes $0\le y-x$, $x<y$ becomes $0\le y-(x+1)$, an equality becomes two facts, a negated inequality is read as the reverse one), and negates the goal.
It searches for integers $m_i\ge0$ and $\kappa\ge1$ such that
\[
\sum_i m_i\,p_i \;\equiv\; -\kappa \qquad\text{as polynomials},
\]
every non-constant monomial cancelling, with a monomial that is not linear treated as an atom.
The proof term is then: $0\le\sum_i m_i p_i$ by the lemmas \texttt{Z.mul\_nonneg} and \texttt{Z.add\_nonneg}; $\sum_i m_ip_i=-\kappa$ by \code{ring}; and the absurdity of $0\le-\kappa$ by \texttt{Z.not\_nonneg\_neg}.
By Farkas' lemma, for a system of affine non-strict inequalities over $\mathbb Q$ such a certificate exists if and only if the system has no solution, so the procedure is complete for such systems.
It is incomplete over $\Z$: the facts $0\le 2x-1$ and $0\le 1-2x$ have no solution in the integers and no certificate, since the procedure has no cutting planes, and the strengthening $x<y\Rightarrow x+1\le y$ recovers only part of integer reasoning.
The search is a linear program (feasibility of $\sum_i m_ip_i=-1$ with $m\ge0$ and a row per non-constant monomial), solved by an integer-preserving simplex method with Bland's rule \citep{bland77} in exact arithmetic; for large problems a floating-point linear-programming solver (HiGHS through SciPy \citep{huangfu18,scipy20}) first guesses the support, which the exact simplex then re-solves, so that floating-point arithmetic decides nothing.

\paragraph{Nonlinear arithmetic.}
\code{nlinarith} adds, before the search, the squares $0\le t^2$ of the subterms $t\cdot t$ and $t^2$ occurring in the hypotheses and the products $0\le p_ip_j$ of every pair of facts (at most $600$), and then searches as above.
The certificate is of the form $\sum_i\lambda_ip_i+\sum_{i\le j}\mu_{ij}p_ip_j+\sum_k\nu_kt_k^2\equiv-\kappa$, a degree-two certificate in the style of the Positivstellensatz \citep{harrison07}.
It is a heuristic, in the sense that the choice of products is a heuristic; any certificate it finds is checked.

\paragraph{Linear combinations of equations.}
\code{linear\_combination} closes an equation $l=r$ from hypotheses $l_i=r_i$ and given coefficients $c_i$ when $l-r=\sum_ic_i(l_i-r_i)$ as a polynomial identity.
The proof applies the congruence lemmas of the ring to obtain $\sum_ic_il_i=\sum_ic_ir_i$, proves $l+\sum_ic_ir_i=r+\sum_ic_il_i$ by \code{ring} (rearranged so that no subtraction is needed), and cancels the common summand.

\subsection{Metatheory: what is and is not established}\label{sec:meta}

We have not proved subject reduction, strong normalisation, consistency or canonicity for the calculus as implemented, and we have not mechanised any part of its metatheory.
For the pure core, a dependent product with an impredicative bottom universe, a universe hierarchy without cumulativity, and the usual conversion, these results are known in the literature on the Calculus of Constructions and its extensions \citep{coquand88cc,coquand90meta,barendregt92}, and for inductive types under strict positivity and structural recursion in that on the Calculus of Inductive Constructions \citep{coquand90ind,paulin93,gimenez94}.
The kernel's calculus differs from those in features whose metatheory is not covered by such a reference, and each of these is a place where the implementation could be wrong in a way that no result we can cite excludes:
\begin{itemize}
\item the index forcing of \cref{sec:ind}, which changes the expected type of a branch;
\item definitional proof irrelevance decided by the unifier and not by an untyped conversion (\cref{sec:prop}), with a depth bound;
\item the strict-positivity check, which does not look in the indices (D2 in \cref{app:deviations}, where a family that it admits yields a proof of $\sy{False}$ with two classical axioms);
\item the shape-only check of a motive (D1);
\item the eager one-layer unfolding of a fixpoint at any argument, against the specification's unfolding at a constructor (\cref{sec:conv});
\item the incompleteness and non-transitivity of conversion on stuck case analyses (\cref{sec:limits}), which can only reject and cannot accept an equation that evaluation does not justify.
\end{itemize}
Verified checkers exist for neighbouring calculi (MetaCoq for the kernel of Rocq, Lean4Lean for that of Lean \citep{sozeau20,carneiro24}); nothing of that kind exists here.
The evidence that the implementation is sound is therefore a test apparatus: a second checker written from the specification, differential and randomised testing against a frozen reference kernel, mutation testing of the kernel and of the tools, sanitiser runs, and a findings log (\cref{sec:testing}).
Such evidence finds bugs and, as the case of D1 and D2 shows, it also marks its own limits: both checkers share one specification, so a flaw in the specification itself is invisible to the comparison.
